\section{Đa agent}

\subsection{Câu hỏi nghiên cứu}

\textit{Khi nhiều Character sống trong cùng một world, interaction có tạo ra emergent behavior không?}

\subsection{Bằng chứng cho Emergence}

\textbf{CÓ} --- bằng chứng mạnh từ:
\begin{itemize}
    \item Stanford Generative Agents (2023): 25 agents, 5 emergent phenomena
    \item CAREB-MAS (2026): Labor specialization, guanxi ethics, clan stratification
\end{itemize}

\subsection{Kết quả Scaling}

\begin{table}[H]
\centering
\caption{Multi-Agent Scaling}
\label{tab:multi-agent-scale}
\begin{tabular}{lll}
\toprule
\textbf{Agent Count} & \textbf{Coordination Overhead} & \textbf{Qualitative Shift} \\
\midrule
5 & Thấp & Không có \\
7 & Trung bình & Không có \\
25 & Trung bình-cao & Emergence bắt đầu \\
100+ & Cao ($>$50\%) & Cấu trúc xã hội mới \\
\bottomrule
\end{tabular}
\end{table}

\subsection{Khuyến nghị Kiến trúc}

\textbf{Hybrid}: Orchestrator + decentralized clusters
\begin{itemize}
    \item Central orchestrator handles cross-cluster communication
    \item Decentralized clusters handle local interactions
    \item Communication protocol: structured messages + natural language
\end{itemize}

% =============================================================================
% 13. CONTEXT ENGINEERING